% Options for packages loaded elsewhere
% Options for packages loaded elsewhere
\PassOptionsToPackage{unicode}{hyperref}
\PassOptionsToPackage{hyphens}{url}
\PassOptionsToPackage{dvipsnames,svgnames,x11names}{xcolor}
%
\documentclass[
  letterpaper,
  DIV=11,
  numbers=noendperiod]{scrreprt}
\usepackage{xcolor}
\usepackage{amsmath,amssymb}
\setcounter{secnumdepth}{5}
\usepackage{iftex}
\ifPDFTeX
  \usepackage[T1]{fontenc}
  \usepackage[utf8]{inputenc}
  \usepackage{textcomp} % provide euro and other symbols
\else % if luatex or xetex
  \usepackage{unicode-math} % this also loads fontspec
  \defaultfontfeatures{Scale=MatchLowercase}
  \defaultfontfeatures[\rmfamily]{Ligatures=TeX,Scale=1}
\fi
\usepackage{lmodern}
\ifPDFTeX\else
  % xetex/luatex font selection
\fi
% Use upquote if available, for straight quotes in verbatim environments
\IfFileExists{upquote.sty}{\usepackage{upquote}}{}
\IfFileExists{microtype.sty}{% use microtype if available
  \usepackage[]{microtype}
  \UseMicrotypeSet[protrusion]{basicmath} % disable protrusion for tt fonts
}{}
\makeatletter
\@ifundefined{KOMAClassName}{% if non-KOMA class
  \IfFileExists{parskip.sty}{%
    \usepackage{parskip}
  }{% else
    \setlength{\parindent}{0pt}
    \setlength{\parskip}{6pt plus 2pt minus 1pt}}
}{% if KOMA class
  \KOMAoptions{parskip=half}}
\makeatother
% Make \paragraph and \subparagraph free-standing
\makeatletter
\ifx\paragraph\undefined\else
  \let\oldparagraph\paragraph
  \renewcommand{\paragraph}{
    \@ifstar
      \xxxParagraphStar
      \xxxParagraphNoStar
  }
  \newcommand{\xxxParagraphStar}[1]{\oldparagraph*{#1}\mbox{}}
  \newcommand{\xxxParagraphNoStar}[1]{\oldparagraph{#1}\mbox{}}
\fi
\ifx\subparagraph\undefined\else
  \let\oldsubparagraph\subparagraph
  \renewcommand{\subparagraph}{
    \@ifstar
      \xxxSubParagraphStar
      \xxxSubParagraphNoStar
  }
  \newcommand{\xxxSubParagraphStar}[1]{\oldsubparagraph*{#1}\mbox{}}
  \newcommand{\xxxSubParagraphNoStar}[1]{\oldsubparagraph{#1}\mbox{}}
\fi
\makeatother


\usepackage{longtable,booktabs,array}
\newcounter{none} % for unnumbered tables
\usepackage{calc} % for calculating minipage widths
% Correct order of tables after \paragraph or \subparagraph
\usepackage{etoolbox}
\makeatletter
\patchcmd\longtable{\par}{\if@noskipsec\mbox{}\fi\par}{}{}
\makeatother
% Allow footnotes in longtable head/foot
\IfFileExists{footnotehyper.sty}{\usepackage{footnotehyper}}{\usepackage{footnote}}
\makesavenoteenv{longtable}
\usepackage{graphicx}
\makeatletter
\newsavebox\pandoc@box
\newcommand*\pandocbounded[1]{% scales image to fit in text height/width
  \sbox\pandoc@box{#1}%
  \Gscale@div\@tempa{\textheight}{\dimexpr\ht\pandoc@box+\dp\pandoc@box\relax}%
  \Gscale@div\@tempb{\linewidth}{\wd\pandoc@box}%
  \ifdim\@tempb\p@<\@tempa\p@\let\@tempa\@tempb\fi% select the smaller of both
  \ifdim\@tempa\p@<\p@\scalebox{\@tempa}{\usebox\pandoc@box}%
  \else\usebox{\pandoc@box}%
  \fi%
}
% Set default figure placement to htbp
\def\fps@figure{htbp}
\makeatother





\setlength{\emergencystretch}{3em} % prevent overfull lines

\providecommand{\tightlist}{%
  \setlength{\itemsep}{0pt}\setlength{\parskip}{0pt}}



 


\KOMAoption{captions}{tableheading}
\makeatletter
\@ifpackageloaded{tcolorbox}{}{\usepackage[skins,breakable]{tcolorbox}}
\@ifpackageloaded{fontawesome5}{}{\usepackage{fontawesome5}}
\definecolor{quarto-callout-color}{HTML}{909090}
\definecolor{quarto-callout-note-color}{HTML}{0758E5}
\definecolor{quarto-callout-important-color}{HTML}{CC1914}
\definecolor{quarto-callout-warning-color}{HTML}{EB9113}
\definecolor{quarto-callout-tip-color}{HTML}{00A047}
\definecolor{quarto-callout-caution-color}{HTML}{FC5300}
\definecolor{quarto-callout-color-frame}{HTML}{acacac}
\definecolor{quarto-callout-note-color-frame}{HTML}{4582ec}
\definecolor{quarto-callout-important-color-frame}{HTML}{d9534f}
\definecolor{quarto-callout-warning-color-frame}{HTML}{f0ad4e}
\definecolor{quarto-callout-tip-color-frame}{HTML}{02b875}
\definecolor{quarto-callout-caution-color-frame}{HTML}{fd7e14}
\makeatother
\makeatletter
\@ifpackageloaded{bookmark}{}{\usepackage{bookmark}}
\makeatother
\makeatletter
\@ifpackageloaded{caption}{}{\usepackage{caption}}
\AtBeginDocument{%
\ifdefined\contentsname
  \renewcommand*\contentsname{Table of contents}
\else
  \newcommand\contentsname{Table of contents}
\fi
\ifdefined\listfigurename
  \renewcommand*\listfigurename{List of Figures}
\else
  \newcommand\listfigurename{List of Figures}
\fi
\ifdefined\listtablename
  \renewcommand*\listtablename{List of Tables}
\else
  \newcommand\listtablename{List of Tables}
\fi
\ifdefined\figurename
  \renewcommand*\figurename{Figure}
\else
  \newcommand\figurename{Figure}
\fi
\ifdefined\tablename
  \renewcommand*\tablename{Table}
\else
  \newcommand\tablename{Table}
\fi
}
\@ifpackageloaded{float}{}{\usepackage{float}}
\floatstyle{ruled}
\@ifundefined{c@chapter}{\newfloat{codelisting}{h}{lop}}{\newfloat{codelisting}{h}{lop}[chapter]}
\floatname{codelisting}{Listing}
\newcommand*\listoflistings{\listof{codelisting}{List of Listings}}
\makeatother
\makeatletter
\makeatother
\makeatletter
\@ifpackageloaded{caption}{}{\usepackage{caption}}
\@ifpackageloaded{subcaption}{}{\usepackage{subcaption}}
\makeatother
\usepackage{bookmark}
\IfFileExists{xurl.sty}{\usepackage{xurl}}{} % add URL line breaks if available
\urlstyle{same}
\hypersetup{
  pdftitle={Computer Networks},
  pdfauthor={K. R. Kaza},
  colorlinks=true,
  linkcolor={blue},
  filecolor={Maroon},
  citecolor={Blue},
  urlcolor={Blue},
  pdfcreator={LaTeX via pandoc}}


\title{Computer Networks}
\usepackage{etoolbox}
\makeatletter
\providecommand{\subtitle}[1]{% add subtitle to \maketitle
  \apptocmd{\@title}{\par {\large #1 \par}}{}{}
}
\makeatother
\subtitle{Course Notes}
\author{K. R. Kaza}
\date{2026-05-01}
\begin{document}
\maketitle

\renewcommand*\contentsname{Table of contents}
{
\setcounter{tocdepth}{2}
\tableofcontents
}

\bookmarksetup{startatroot}

\chapter*{Preface}\label{preface}
\addcontentsline{toc}{chapter}{Preface}

\markboth{Preface}{Preface}

This is a Quarto book.

To learn more about Quarto books visit
\url{https://quarto.org/docs/books}.

\bookmarksetup{startatroot}

\chapter{Introduction}\label{introduction}

This chapter introduces the basic vocabulary of data communications and
computer networks. We answer four questions: \emph{what does it mean for
two devices to communicate?}, \emph{what is a network?}, \emph{what is
the Internet?}, and \emph{what are protocols and standards, and why do
we need them?} The concepts here will return in every later chapter ---
once you have the vocabulary, everything that follows is filling in
details.

\section{Data communications}\label{sec-data-comm}

\textbf{Data communications} is the exchange of data between two devices
over some form of transmission medium such as a wire, optical fibre, or
radio link. For communication to occur, three conditions must be met:
the devices must be part of a communication system made up of a
combination of hardware (physical equipment) and software (programs),
the data must be delivered to the correct destination, and the data must
be delivered accurately and in a timely manner.

\subsection{Components of a communication
system}\label{components-of-a-communication-system}

Any data communication system has five essential components, illustrated
in Figure~\ref{fig-components}.

\begin{figure}

\centering{

\includegraphics[width=8.29in,height=1.75in]{intro_files/figure-latex/mermaid-figure-1.png}

}

\caption{\label{fig-components}Five components of a communication
system.}

\end{figure}%

\begin{enumerate}
\def\labelenumi{\arabic{enumi}.}
\tightlist
\item
  \textbf{Message} --- the information to be communicated. It can be
  text, numbers, images, audio, or video.
\item
  \textbf{Sender} --- the device that sends the message (a computer,
  phone, camera, sensor, etc.).
\item
  \textbf{Receiver} --- the device that receives the message.
\item
  \textbf{Transmission medium} --- the physical path over which the
  message travels (twisted-pair cable, coaxial cable, optical fibre,
  radio waves).
\item
  \textbf{Protocol} --- a set of rules that governs the communication.
  Without a shared protocol, two devices may be connected but cannot
  understand each other.
\end{enumerate}

\begin{tcolorbox}[enhanced jigsaw, arc=.35mm, bottomrule=.15mm, bottomtitle=1mm, breakable, colback=white, colbacktitle=quarto-callout-note-color!10!white, colframe=quarto-callout-note-color-frame, coltitle=black, left=2mm, leftrule=.75mm, opacityback=0, opacitybacktitle=0.6, rightrule=.15mm, title=\textcolor{quarto-callout-note-color}{\faInfo}\hspace{0.5em}{Note}, titlerule=0mm, toprule=.15mm, toptitle=1mm]

A person speaking French to a person who only understands English
illustrates the need for a protocol: the medium (air) carries the
message, the sender produces it, the receiver hears it, but
communication fails because the rules of encoding are different.

\end{tcolorbox}

\subsection{Data representation}\label{data-representation}

Inside a digital system, every message --- regardless of its original
form --- is represented as a sequence of bits.

\begin{itemize}
\tightlist
\item
  \textbf{Text} is encoded using character sets such as ASCII (7 bits
  per character) or Unicode (typically 8, 16, or 32 bits per character
  via UTF-8/16/32).
\item
  \textbf{Numbers} are encoded directly in binary.
\item
  \textbf{Images} are represented as a matrix of pixels; each pixel is
  one or more bytes describing its colour.
\item
  \textbf{Audio} is sampled at a fixed rate (e.g.~44.1 kHz for
  CD-quality) and each sample is quantised to a fixed bit depth.
\item
  \textbf{Video} is a sequence of images (frames) shown rapidly, usually
  compressed using a codec such as H.264 or AV1.
\end{itemize}

\subsection{Data flow}\label{data-flow}

Communication between two devices can be one of three types, summarised
in Table~\ref{tbl-flow}.

\begin{longtable}[]{@{}
  >{\raggedright\arraybackslash}p{(\linewidth - 4\tabcolsep) * \real{0.1807}}
  >{\raggedright\arraybackslash}p{(\linewidth - 4\tabcolsep) * \real{0.5301}}
  >{\raggedright\arraybackslash}p{(\linewidth - 4\tabcolsep) * \real{0.2892}}@{}}
\caption{Data flow modes.}\label{tbl-flow}\tabularnewline
\toprule\noalign{}
\begin{minipage}[b]{\linewidth}\raggedright
Mode
\end{minipage} & \begin{minipage}[b]{\linewidth}\raggedright
Direction
\end{minipage} & \begin{minipage}[b]{\linewidth}\raggedright
Example
\end{minipage} \\
\midrule\noalign{}
\endfirsthead
\toprule\noalign{}
\begin{minipage}[b]{\linewidth}\raggedright
Mode
\end{minipage} & \begin{minipage}[b]{\linewidth}\raggedright
Direction
\end{minipage} & \begin{minipage}[b]{\linewidth}\raggedright
Example
\end{minipage} \\
\midrule\noalign{}
\endhead
\bottomrule\noalign{}
\endlastfoot
Simplex & One direction only & Keyboard → computer \\
Half-duplex & Both directions, but only one at a time & Walkie-talkie \\
Full-duplex & Both directions simultaneously & Telephone call \\
\end{longtable}

\section{Networks}\label{sec-networks}

A \textbf{network} is a set of devices (often called \emph{nodes})
connected by communication links. A node can be a computer, printer,
phone, or any other device capable of sending or receiving data on the
network.

\subsection{Network criteria}\label{network-criteria}

A network is judged by how well it meets four criteria.

\begin{description}
\tightlist
\item[Performance]
Measured by \textbf{throughput} (how much data crosses the network per
unit time) and \textbf{delay} (how long a unit of data takes to cross).
These two metrics are often in tension --- increasing throughput by
sending more data may also increase delay because of congestion.
\item[Reliability]
Measured by the frequency of failures, the time it takes to recover from
a failure, and the network's robustness in a catastrophe.
\item[Security]
The ability to protect data from unauthorised access and damage, and to
recover from data loss.
\item[Scalability]
The ability of the network to grow --- to add more nodes or carry more
traffic --- without redesign.
\end{description}

\subsection{Physical structures}\label{physical-structures}

The physical structure of a network has two aspects: \emph{how} devices
are connected and \emph{what shape} the connections take.

\textbf{Type of connection.} A link can be either:

\begin{itemize}
\tightlist
\item
  \emph{Point-to-point}: a dedicated link between exactly two devices
  (e.g.~the cable from your laptop to its monitor).
\item
  \emph{Multipoint}: a single link shared by more than two devices
  (e.g.~classical Ethernet on a coaxial bus).
\end{itemize}

\textbf{Physical topology.} The geometric layout of the links. The four
primary topologies are:

\begin{itemize}
\tightlist
\item
  \textbf{Mesh} --- every device has a dedicated link to every other
  device. For \(n\) nodes this requires \(\dfrac{n(n-1)}{2}\) links.
  Extremely robust and fast, but expensive and unscalable.
\item
  \textbf{Star} --- every device has a dedicated link to a central hub.
  The hub becomes a single point of failure but the topology is cheap
  and easy to manage.
\item
  \textbf{Bus} --- all devices share a single backbone cable. Cheap and
  simple, but the backbone is a single point of failure.
\item
  \textbf{Ring} --- each device is connected to two neighbours, forming
  a closed loop. Signals travel in one direction around the ring.
\item
  \textbf{Hybrid} --- a mix of the above; for example, a star of buses.
\end{itemize}

\begin{tcolorbox}[enhanced jigsaw, arc=.35mm, bottomrule=.15mm, bottomtitle=1mm, breakable, colback=white, colbacktitle=quarto-callout-tip-color!10!white, colframe=quarto-callout-tip-color-frame, coltitle=black, left=2mm, leftrule=.75mm, opacityback=0, opacitybacktitle=0.6, rightrule=.15mm, title=\textcolor{quarto-callout-tip-color}{\faLightbulb}\hspace{0.5em}{Counting links in a mesh}, titlerule=0mm, toprule=.15mm, toptitle=1mm]

A common exam question: \emph{if a mesh network has \(n=10\) nodes, how
many links are required?} Plug into the formula:
\(10 \times 9 / 2 = 45\) links, and each node needs \(n-1 = 9\) I/O
ports. Mesh becomes impractical fast --- this is why almost no real
network is fully meshed.

\end{tcolorbox}

\subsection{Network types by reach}\label{network-types-by-reach}

Networks are classified by the geographic area they cover.

\begin{longtable}[]{@{}
  >{\raggedright\arraybackslash}p{(\linewidth - 4\tabcolsep) * \real{0.0833}}
  >{\raggedright\arraybackslash}p{(\linewidth - 4\tabcolsep) * \real{0.2778}}
  >{\raggedright\arraybackslash}p{(\linewidth - 4\tabcolsep) * \real{0.6389}}@{}}
\caption{Networks classified by reach.}\label{tbl-reach}\tabularnewline
\toprule\noalign{}
\begin{minipage}[b]{\linewidth}\raggedright
Type
\end{minipage} & \begin{minipage}[b]{\linewidth}\raggedright
Reach
\end{minipage} & \begin{minipage}[b]{\linewidth}\raggedright
Typical use
\end{minipage} \\
\midrule\noalign{}
\endfirsthead
\toprule\noalign{}
\begin{minipage}[b]{\linewidth}\raggedright
Type
\end{minipage} & \begin{minipage}[b]{\linewidth}\raggedright
Reach
\end{minipage} & \begin{minipage}[b]{\linewidth}\raggedright
Typical use
\end{minipage} \\
\midrule\noalign{}
\endhead
\bottomrule\noalign{}
\endlastfoot
PAN & A few metres & Bluetooth headset, fitness tracker \\
LAN & A building, campus & Office Ethernet, home Wi-Fi \\
MAN & A city & City-wide Wi-Fi, cable TV distribution \\
WAN & A country or globe & The public Internet, corporate backbone \\
\end{longtable}

A \textbf{LAN} is typically owned and operated by a single organisation.
A \textbf{WAN} crosses public infrastructure and is usually built from
links leased from telecommunications providers. The Internet is a
network of LANs and WANs glued together.

\subsection{Switching}\label{sec-switching}

A network of more than a few devices cannot afford a dedicated link
between every pair of nodes. Instead, intermediate devices called
\emph{switches} relay data along a path from source to destination.
There are two fundamental strategies:

\begin{itemize}
\tightlist
\item
  \textbf{Circuit switching} --- before any data is sent, a dedicated
  path of links is reserved end-to-end for the duration of the
  conversation, and a fixed share of each link's capacity is set aside
  for it. This is how traditional telephone networks worked.
\item
  \textbf{Packet switching} --- the message is broken into small chunks
  called \emph{packets}. Each packet is forwarded independently through
  the network. Packets from many conversations share the same links
  statistically.
\end{itemize}

The Internet is a packet-switched network. The advantage is high
utilisation: idle moments in one conversation are filled by traffic from
another. The disadvantage is variability: when many packets converge on
a router, some have to wait in a queue, and if the queue overflows,
packets are \emph{lost}.

\section{The Internet}\label{sec-internet}

The \textbf{Internet} is a global packet-switched network of networks
that uses the TCP/IP protocol suite. It is not a single network owned by
anyone; it is the union of many independent networks that agree to
interconnect.

\subsection{A very brief history}\label{a-very-brief-history}

\begin{itemize}
\tightlist
\item
  \textbf{1960s} --- ARPANET, funded by the U.S. Defense Department,
  demonstrates that packet switching works.
\item
  \textbf{1970s} --- Vint Cerf and Bob Kahn design TCP/IP, a protocol
  that allows different networks to interconnect.
\item
  \textbf{1983} --- ARPANET officially switches to TCP/IP. This date is
  often cited as the birth of the Internet.
\item
  \textbf{1989} --- Tim Berners-Lee proposes the World Wide Web at CERN.
\item
  \textbf{1990s} --- The Web, e-mail, and commercial ISPs bring the
  Internet to the general public.
\item
  \textbf{2000s--today} --- Mobile, cloud computing, video streaming,
  and the Internet of Things become dominant traffic sources.
\end{itemize}

\subsection{The Internet today}\label{the-internet-today}

The modern Internet has a hierarchical structure of \emph{Internet
Service Providers} (ISPs).

\begin{itemize}
\tightlist
\item
  \textbf{End users} (homes, businesses, mobile devices) connect to a
  \textbf{local ISP}.
\item
  Local ISPs connect to \textbf{regional ISPs}.
\item
  Regional ISPs connect to \textbf{Tier-1 ISPs} (also called
  \emph{backbone} providers), which span continents and exchange traffic
  at \emph{Internet exchange points} (IXPs).
\end{itemize}

This hierarchy is not strict --- large content providers (Google,
Netflix, Meta) operate their own private backbones and peer directly
with many ISPs, bypassing the upper tiers. But the mental model of
\emph{edge → access → core} is still useful.

\section{Protocols and standards}\label{sec-protocols}

We have used the word \emph{protocol} informally; let us define it.

A \textbf{protocol} is a set of rules that governs data communication.
Every protocol defines three things:

\begin{description}
\tightlist
\item[Syntax]
The structure or format of the data --- the order in which bits are
transmitted, the meaning of each field, the size of headers and
payloads.
\item[Semantics]
The meaning of each piece of the message. For example, an address field
tells the receiver \emph{who is being addressed}; an error-checking
field tells the receiver \emph{whether the data was corrupted in
transit}.
\item[Timing]
When data should be sent and how fast. If a sender produces 100 Mbps but
the receiver can only consume 1 Mbps, the data will be lost without
timing rules to slow the sender down.
\end{description}

\subsection{Standards}\label{standards}

A \textbf{standard} is a protocol (or set of protocols) that has been
formally agreed upon. Standards are essential --- without them,
equipment from different vendors could not interoperate.

Standards come in two flavours:

\begin{itemize}
\tightlist
\item
  \textbf{De facto} standards arise from widespread use, without any
  formal endorsement. The QWERTY keyboard layout is a de facto standard.
\item
  \textbf{De jure} standards are formally adopted by a recognised body.
  The IEEE 802.11 family of Wi-Fi standards is de jure.
\end{itemize}

The main standards organisations relevant to networking are summarised
in Table~\ref{tbl-standards}.

\begin{longtable}[]{@{}
  >{\raggedright\arraybackslash}p{(\linewidth - 4\tabcolsep) * \real{0.0642}}
  >{\raggedright\arraybackslash}p{(\linewidth - 4\tabcolsep) * \real{0.5046}}
  >{\raggedright\arraybackslash}p{(\linewidth - 4\tabcolsep) * \real{0.4312}}@{}}
\caption{Major standards bodies in
networking.}\label{tbl-standards}\tabularnewline
\toprule\noalign{}
\begin{minipage}[b]{\linewidth}\raggedright
Body
\end{minipage} & \begin{minipage}[b]{\linewidth}\raggedright
Stands for
\end{minipage} & \begin{minipage}[b]{\linewidth}\raggedright
What it standardises
\end{minipage} \\
\midrule\noalign{}
\endfirsthead
\toprule\noalign{}
\begin{minipage}[b]{\linewidth}\raggedright
Body
\end{minipage} & \begin{minipage}[b]{\linewidth}\raggedright
Stands for
\end{minipage} & \begin{minipage}[b]{\linewidth}\raggedright
What it standardises
\end{minipage} \\
\midrule\noalign{}
\endhead
\bottomrule\noalign{}
\endlastfoot
ISO & International Organization for Standardization & The OSI reference
model, many others \\
ITU-T & International Telecommunication Union --- Telecom & Telephony,
optical transport \\
IEEE & Institute of Electrical and Electronics Engineers & LAN/MAN
standards: Ethernet (802.3), Wi-Fi (802.11) \\
IETF & Internet Engineering Task Force & Internet protocols, published
as RFCs \\
ANSI & American National Standards Institute & Many U.S. standards \\
\end{longtable}

Internet standards are published as \textbf{Requests for Comments}
(RFCs). Despite the modest name, an RFC that reaches \emph{Internet
Standard} status is binding for anyone who wants to interoperate. For
example, RFC 791 defines IPv4; RFC 9293 defines TCP.

\section{Summary}\label{summary}

\begin{itemize}
\tightlist
\item
  \textbf{Data communication} requires five components: a message, a
  sender, a receiver, a transmission medium, and a protocol.
\item
  A \textbf{network} is a set of devices connected by communication
  links. It is judged on performance, reliability, security, and
  scalability.
\item
  Networks differ in \textbf{topology} (mesh, star, bus, ring, hybrid)
  and in \textbf{reach} (PAN, LAN, MAN, WAN).
\item
  The Internet is a global \textbf{packet-switched} network of networks,
  organised as a loose hierarchy of ISPs.
\item
  A \textbf{protocol} specifies syntax, semantics, and timing; a
  \textbf{standard} is a protocol that everyone has agreed to follow.
\end{itemize}

The remaining chapters look at each layer of the protocol stack in turn,
starting from the bottom --- the physical layer --- and working upward.

\section*{Exercises}\label{exercises}
\addcontentsline{toc}{section}{Exercises}

\markright{Exercises}

\begin{enumerate}
\def\labelenumi{\arabic{enumi}.}
\tightlist
\item
  List the five components of a data communication system and explain
  the role of each.
\item
  A full mesh network has 12 nodes. How many links does it use? How many
  I/O ports does each node need?
\item
  A camera streams 1080p video at 30 frames per second. Each frame is
  \(1920 \times 1080\) pixels, and each pixel uses 24 bits of colour.
  What raw bit rate does this produce, in Mbps? Why is the actual
  transmitted bit rate much lower?
\item
  Classify each of the following as simplex, half-duplex, or
  full-duplex: (a) a TV broadcast, (b) two children with a tin-can
  telephone, (c) a video call, (d) a remote control and a TV, (e) a CB
  radio conversation.
\item
  The Internet is described as a ``network of networks''. What does this
  mean? Why is no single organisation in charge of it?
\end{enumerate}




\end{document}
